% !TEX program = pdflatex
\documentclass[11pt,a4paper]{article}

% =========================================================================
% ACTA DE EVALUACIÓN DE CARTERA
% Documento común a los proyectos CasaScore y Horus Sentinel
% Gestión de Proyectos TICS (EST315) -- Semestre 2026-II -- Team Sobrevivientes
% =========================================================================

% --- PAQUETES DE CONFIGURACIÓN ---
\usepackage[utf8]{inputenc}
\usepackage[T1]{fontenc}
\usepackage[spanish,es-tabla]{babel}
% Compatibilidad con babel-spanish reciente (flechas de TikZ y signo de porcentaje)
\AtBeginDocument{\shorthandoff{<>}\spanishplainpercent}
\usepackage{geometry}
\geometry{
    top=3.6cm,
    bottom=2.3cm,
    left=2.2cm,
    right=2.2cm,
    % PARTE MODIFICABLE: headheight debe ser mayor que la altura del logo más alto (2.56cm = 73pt).
    headheight=80pt,
    headsep=15pt,
    footskip=20pt
}

\usepackage{graphicx}
\usepackage{fancyhdr}
\usepackage[table]{xcolor}
\usepackage{titlesec}
\usepackage{booktabs}
\usepackage{tabularx}
\usepackage{array}
\usepackage{amsmath}
\usepackage{enumitem}
\usepackage{microtype}
\usepackage{hyperref}

% Logos: funciona a 2 o 3 niveles de profundidad dentro de docs/latex/
\graphicspath{{../../assets/img/}{../../../assets/img/}}

% --- PALETA INSTITUCIONAL FIEI / UNA PUNO ---
\definecolor{AzulUNA}{RGB}{10, 45, 90}
\definecolor{DoradoFIEI}{RGB}{180, 130, 20}
\definecolor{GrisOscuro}{RGB}{50, 50, 50}
\definecolor{AzulClaro}{RGB}{228, 236, 246}

% --- CONFIGURACIÓN DE LISTAS ---
\setlist{topsep=2pt, itemsep=2pt, parsep=1pt}

% --- COLUMNAS DE TABLA ---
\newcolumntype{L}{>{\raggedright\arraybackslash}X}
\newcolumntype{P}[1]{>{\raggedright\arraybackslash}p{#1}}
\newcolumntype{C}[1]{>{\centering\arraybackslash}p{#1}}
\renewcommand{\arraystretch}{1.25}

% --- FORMATO DE ENLACES ---
\hypersetup{
    colorlinks=true,
    linkcolor=AzulUNA,
    urlcolor=AzulUNA,
    citecolor=AzulUNA,
    pdftitle={Acta de Evaluación de Cartera},
    pdfauthor={Team Sobrevivientes}
}

% --- FORMATO DE TÍTULOS ---
\titleformat{\section}{\color{AzulUNA}\normalfont\large\bfseries}{\thesection.}{0.7em}{}[\color{AzulUNA}\titlerule]
\titleformat{\subsection}{\color{AzulUNA}\normalfont\normalsize\bfseries}{\thesubsection.}{0.5em}{}
\titlespacing*{\section}{0pt}{9pt}{4pt}
\titlespacing*{\subsection}{0pt}{6pt}{2pt}

% --- ENCABEZADO Y PIE DE PÁGINA (fancyhdr) ---
\pagestyle{fancy}
\fancyhf{}

% ===== TAMAÑO DE LOS LOGOS (PARTE MODIFICABLE) =====
% logo_fiei.png  : height=1.45cm (lado izquierdo)
% logo_equipo.png: height=2.56cm (lado derecho; 25 % más grande porque la imagen tiene más borde blanco)
% Para cambiar el tamaño, modifique solo el valor height=... de cada \includegraphics.
% Si el logo más alto supera 2.6cm, aumente también headheight y top en la configuración de página.
\fancyhead[L]{%
    \includegraphics[height=1.45cm,keepaspectratio]{logo_fiei.png}%
}
\fancyhead[C]{%
    \scriptsize \textbf{\textcolor{AzulUNA}{UNIVERSIDAD NACIONAL DEL ALTIPLANO DE PUNO}}\\[2pt]
    \tiny \textcolor{GrisOscuro}{FACULTAD DE INGENIERÍA ESTADÍSTICA E INFORMÁTICA}\\[2pt]
    \scriptsize \textbf{\textcolor{AzulUNA}{Gestión de Proyectos TICS (EST315) \textbar\ Team Sobrevivientes}}\\[2pt]
    \footnotesize \textbf{\textcolor{DoradoFIEI}{ACTA DE EVALUACIÓN DE CARTERA}}%
}
\fancyhead[R]{%
    \includegraphics[height=2.56cm,keepaspectratio]{logo_equipo.png}%
}

\renewcommand{\headrulewidth}{0pt}
\renewcommand{\headrule}{}

\fancyfoot[L]{\scriptsize \textcolor{GrisOscuro}{CasaScore \textperiodcentered\ Horus Sentinel}}
\fancyfoot[C]{\small \thepage}
\fancyfoot[R]{\scriptsize \textcolor{GrisOscuro}{Puno, Perú -- Semestre 2026-II}}
\renewcommand{\footrulewidth}{0.4pt}

\begin{document}

% =========================================================================
% CABECERA DEL DOCUMENTO
% =========================================================================

\begin{center}
    {\Large \textbf{\textcolor{AzulUNA}{ACTA DE EVALUACIÓN DE CARTERA}}}\\[3pt]
    {\normalsize \textbf{Evaluación de las propuestas CasaScore y Horus Sentinel y decisión de continuidad}}\\[4pt]
    \footnotesize \textbf{Fecha:} martes 08/09/2026 (semana 3) \quad \textbar\quad \textbf{Hora:} 20:00 -- 21:00 \quad \textbar\quad \textbf{Decisión:} aprobada por unanimidad
\end{center}

\vspace{0.1cm}

\noindent\begin{tabularx}{\textwidth}{@{}P{3.8cm} L@{}}
\toprule
\textbf{Curso:} & Gestión de Proyectos TICS (EST315) -- Semestre 2026-II \\
\textbf{Docente:} & Ing. Renzo Apaza Cutipa \\
\textbf{Equipo:} & Team Sobrevivientes (los cuatro integrantes presentes) \\
\textbf{Insumos evaluados:} & Fichas técnicas de ambas propuestas y especificación de requerimientos (04/09/2026). \\
\bottomrule
\end{tabularx}

% =========================================================================
\section{Antecedentes}
% =========================================================================
El equipo formuló dos propuestas de proyecto TIC:
\begin{itemize}[leftmargin=*]
    \item \textbf{CasaScore (AlquilaPuno):} recomendación multicriterio de cuartos en alquiler para estudiantes de la UNA Puno, con puntaje explicado de 0 a 100.
    \item \textbf{Horus Sentinel:} mapa de calor de riesgo delictivo por zona y horario, con sugerencia de rutas seguras.
\end{itemize}
Con ambas fichas y sus requerimientos, el equipo evaluó la viabilidad de cada propuesta frente al plazo del semestre (semanas 2 a 17), la disponibilidad de datos, la complejidad técnica y el costo.

% =========================================================================
\section{Matriz de Evaluación Ponderada}
% =========================================================================
Puntaje de 1 (muy bajo) a 5 (muy alto). Para costo y complejidad, 5 significa más favorable.

\vspace{0.1cm}
{\small
\noindent\begin{tabularx}{\textwidth}{@{}L C{1.1cm} C{1.55cm} C{1.85cm} C{1.55cm} C{1.85cm}@{}}
\toprule
\rowcolor{AzulClaro}
& & \multicolumn{2}{c}{\textbf{CasaScore}} & \multicolumn{2}{c}{\textbf{Horus Sentinel}} \\
\rowcolor{AzulClaro}
\textbf{Criterio} & \textbf{Peso} & \textbf{Puntaje} & \textbf{Ponderado} & \textbf{Puntaje} & \textbf{Ponderado} \\
\midrule
Impacto social & 25\,\% & 4 & 1.00 & 5 & 1.25 \\
Viabilidad técnica (complejidad) & 20\,\% & 5 & 1.00 & 3 & 0.60 \\
Disponibilidad de datos & 20\,\% & 5 & 1.00 & 2 & 0.40 \\
Factibilidad en el semestre & 20\,\% & 5 & 1.00 & 3 & 0.60 \\
Costo de infraestructura & 15\,\% & 5 & 0.75 & 4 & 0.60 \\
\midrule
\textbf{Total} & \textbf{100\,\%} & & \textbf{4.75} & & \textbf{3.45} \\
\bottomrule
\end{tabularx}
}

\vspace{0.1cm}
\noindent\textbf{Regla de decisión:} 4.0 o más, viable; de 3.0 a 3.9, viable con ajustes de alcance; menos de 3.0, no viable.

% =========================================================================
\section{Análisis de Cada Propuesta}
% =========================================================================

\subsection{CasaScore}
\begin{itemize}[leftmargin=*]
    \item \textbf{Datos:} se obtienen directamente en campo, con propietarios y estudiantes de Puno.
    \item \textbf{Modelo:} normalización 0--100 y suma ponderada; es controlable y verificable.
    \item \textbf{Plazo:} un MVP es alcanzable en 5 sprints semanales (semanas 4 a 8).
    \item \textbf{Costo:} servicios en capa gratuita.
\end{itemize}

\subsection{Horus Sentinel}
{\small
\noindent\begin{tabularx}{\textwidth}{@{}P{3.8cm} L L@{}}
\toprule
\rowcolor{AzulClaro}
\textbf{Dificultad} & \textbf{Riesgo} & \textbf{Ajuste aprobado} \\
\midrule
Acceso a datos de criminalidad & Los datos georreferenciados de la Policía son reservados y gestionar un convenio excede el semestre. & Unidad I con datos simulados; MVP final con \textbf{reportes ciudadanos} registrados en la propia aplicación. \\
\addlinespace[2pt]
Complejidad del modelo & El modelo espaciotemporal completo requiere meses de calibración. & Mapa de calor por densidad (KDE) por zona y horario; ruta segura con $A^*$ y costo $= \alpha \cdot \text{distancia} + \beta \cdot \text{riesgo}$. El modelo espaciotemporal avanzado queda fuera del alcance. \\
\addlinespace[2pt]
Sesgo de los datos & Los reportes pueden concentrarse en ciertas zonas o ser falsos. & Reportes anónimos y moderados, límite de reportes por usuario y auditoría periódica del mapa. \\
\addlinespace[2pt]
Cobertura geográfica & Cubrir toda la ciudad eleva el costo de cómputo. & Distrito piloto: \textbf{Puno (cercado)}. \\
\bottomrule
\end{tabularx}
}

% =========================================================================
\section{Resoluciones}
% =========================================================================
\begin{enumerate}[leftmargin=*]
    \item \textbf{Continuidad de ambos proyectos:} CasaScore y Horus Sentinel continúan como dos proyectos separados, cada uno con su propio expediente, línea base y registro de riesgos.
    \item \textbf{Prioridad:} \textbf{CasaScore es el proyecto prioritario} (puntaje 4.75). Cuando ambos compitan por el tiempo del equipo, se atiende primero a CasaScore.
    \item \textbf{Alcance ajustado de Horus Sentinel} (puntaje 3.45, viable con ajustes):
    \begin{itemize}
        \item \textbf{Unidad I (hasta la semana 8):} prototipo funcional con datos simulados del distrito de Puno (cercado).
        \item \textbf{Cierre del semestre (semana 17):} MVP funcional con mapa de calor y ruta segura alimentados por reportes ciudadanos.
    \end{itemize}
    \item \textbf{Dedicación:} cada integrante dedica 8 horas semanales a cada proyecto.
    \item \textbf{Costo:} ambos proyectos usan servicios en capa gratuita (mapas abiertos y base de datos en la nube).
    \item \textbf{Registro:} los acuerdos se publican en el repositorio y se crean los Issues de inicio de cada proyecto.
\end{enumerate}

% =========================================================================
\section{Firmas}
% =========================================================================
\noindent En señal de conformidad con las resoluciones adoptadas, los integrantes del equipo firman la presente acta:

\vspace{1.5cm}

\begin{center}
\begin{tabular}{c c c}
    \parbox{6.5cm}{\centering \rule{6.0cm}{0.6pt}\\[5pt] \textbf{Clyde Neil Paricahua Pari}\\[2pt] {\footnotesize Project Manager}}
    & \hspace{0.8cm} &
    \parbox{6.5cm}{\centering \rule{6.0cm}{0.6pt}\\[5pt] \textbf{Ruth Karina Apaza Solis}\\[2pt] {\footnotesize Analista de Planificación}}
    \\[2.0cm]
    \parbox{6.5cm}{\centering \rule{6.0cm}{0.6pt}\\[5pt] \textbf{Eliseo Tarqui Ajahuana}\\[2pt] {\footnotesize Especialista Técnico / Arquitectura}}
    & \hspace{0.8cm} &
    \parbox{6.5cm}{\centering \rule{6.0cm}{0.6pt}\\[5pt] \textbf{Viviana Mary Quisbert Quispe}\\[2pt] {\footnotesize Gestora de Comunicación, Riesgos y Documentación}}
\end{tabular}
\end{center}

\end{document}
